\begin{figure}[H]
\centering
\caption{Matriz de confusão para uma classe examinada.}
\label{fig:matriz-confusao}
\setlength{\unitlength}{1mm}
\resizebox{0.7\textwidth}{!}{%
\begin{picture}(115,64)
\small
\put(35,52){\makebox(35,12){\shortstack{Classe prevista:\\positiva}}}
\put(70,52){\makebox(35,12){\shortstack{Classe prevista:\\negativa}}}
\put(0,30){\makebox(32,22){\shortstack{Classe\\observada:\\positiva}}}
\put(0,8){\makebox(32,22){\shortstack{Classe\\observada:\\negativa}}}
\put(35,30){\framebox(35,22){\shortstack{VP\\(verdadeiro positivo)}}}
\put(70,30){\framebox(35,22){\shortstack{FN\\(falso negativo)}}}
\put(35,8){\framebox(35,22){\shortstack{FP\\(falso positivo)}}}
\put(70,8){\framebox(35,22){\shortstack{VN\\(verdadeiro negativo)}}}
\put(0,0){\makebox(115,7){\scriptsize Precisão: VP/(VP+FP). Revocação: VP/(VP+FN).}}
\end{picture}
}%
\fonte{Elaborado pelo autor (2026).}
\end{figure}
